% Generated from locale/tr/content/history/biographies/bertrand-russell.tex; do not hand-edit.


\olfileid[tr]{his}{bio}{rus} 

\olsection{Bertrand Russell}

\olphoto{russell-bertrand}{Bertrand Russell}

Bertrand Russell, modern analitik felsefenin kurucularından biri olarak
anılır. 18 Mayıs 1872'de doğan Russell yalnızca felsefe ve mantık alanındaki
çalışmalarıyla tanınmakla kalmadı, çeşitli konularda halka yönelik birçok
kitap da yazdı. Yaşamı boyunca ateşli bir siyasal eylemciydi.

Russell Galler'in Monmouthshire bölgesindeki Trellech'te doğdu. Anne babası
Britanya soyluları sınıfındandı. Özgür düşünceliydiler; hatta dönemin
Boston'ındaki radikallerle dostluk kurmuşlardı. Ne yazık ki Russell küçükken
anne babası öldü ve büyükanne ile büyükbabasının yanında yaşamaya gönderildi.
Orada dinsel bir eğitim gördü (anne babasının ne pahasına olursa olsun
kaçınmak istediği bir şeydi bu). Büyükannesi ahlakla ilgili her konuda çok
katıydı. Ergenlik yıllarında çoğunlukla özel öğretmenlerce evde eğitildi.

Russell'ın analitik felsefe, özellikle de mantık üzerindeki etkisi çok
büyüktür. Cambridge'deki Trinity College'da matematik ve felsefe okudu; burada
matematikçi ve filozof Alfred North Whitehead'den etkilendi. Russell ile
Whitehead 1910'da, matematiğin mantığa indirgenebileceği görüşünü savundukları
\emph{Principia Mathematica}'nın ilk cildini yayımladı. Russell daha sonra
yüzlerce kitap, deneme ve siyasal broşür yayımladı. 1950'de Nobel Edebiyat
Ödülü'nü kazandı.

Russell siyaset ve toplumsal eylemciliğin derinden içindeydi. I. Dünya Savaşı
sırasında pasifist eylemleri ve protestoları nedeniyle tutuklandı ve altı ay
hapse gönderildi. Hapisteyken okuyup yazabildi ve bu deneyimi ``oldukça hoş''
bulduğunu söyledi. Yaşamı boyunca pasifist kalmayı sürdürdü; 1961'de nükleer
silahsızlanma mitingine katıldığı için yeniden hapsedildi. 1948'de, yalnızca
sigara içilen bölümde oturanların kurtulduğu bir uçak kazasından da sağ çıktı.
Bu yüzden Russell hayatını sigaraya borçlu olduğunu ileri sürdü. Dört kez
evlendi; ancak evlilik dışı ilişkiler sürdürmesiyle de ün saldı. 2 Şubat
1970'te Galler'in Penrhyndeudraeth kasabasında 97 yaşında öldü.

\begin{reading}
Russell, yaşamının 1872--1967 yıllarını kapsayan üç bölümlük bir özyaşam
öyküsü yazdı \citep{Russell1967,Russell1968,Russell1969}. McMaster
Üniversitesi'ndeki Bertrand Russell Araştırma Merkezi, Bertrand Russell
arşivlerine ev sahipliği yapar. Toplu eserlerinin ciltleri (aranabilir
dizinler dâhil) ve arşiv projeleri hakkında bilgi için merkezin
\citet{Duncan2015} adresindeki internet sitesine bakın. Russell'ın
\emph{On Denoting} başlıklı makalesi \citep{Russell1905}, 20. yüzyıl analitik
felsefesinin klasiklerindendir.

Stanford Felsefe Ansiklopedisi'nin Russell maddesinde \citep{Irvine2015},
Russell'ın Arzu ve Siyasal Teori üzerine konuşmalarından ses kayıtları
bulunur. Russell'la yapılmış birçok görüntülü söyleşi çevrimiçi olarak
mevcuttur. Örneğin sigara içmekten ve bir uçak kazasına karışmaktan söz
edişini görmek için \citet{RussellND} kaynağına bakın. Russell'ın
\emph{Introduction to Mathematical Philosophy} dâhil bazı eserleri,
\citet{LibriVoxND} üzerinde ücretsiz sesli kitap olarak mevcuttur.
\end{reading}

